\documentclass[a4paper, 12pt]{article}

\usepackage{utils}

\renewcommand*{\today}{25 March 2026}

\begin{document}

\hotbox{Analyse 4}{CM 12}{\today}

\subsection{Approximation $L^2$}

\begin{definition}
    On définit le produit
    $$
    \forall f, g \in Cpm([0, 2\pi[, \mathbb{C}), \quad \langle f, g \rangle = \frac{1}{2\pi} \int_0^{2\pi} f(t) \overline{g(t)} \, dt.
    $$

    On peut définir une norme (dite norme $L^2$) par
    $$
    \norm{f}_2 = \sqrt{\langle f, f \rangle} = \sqrt{\frac{1}{2\pi} \int_0^{2\pi} |f(t)|^2 \, dt}.
    $$
\end{definition}

\begin{lemme}{Inégalité de Schwarz}{}
    Soient $f, g \in Cpm([0, 2\pi[, \mathbb{C})$ telles que $\norm{f}_2 \neq 0$ et $\norm{g}_2 \neq 0$. Alors
    $$
    |\langle f, g \rangle| \leq \norm{f}_2 \norm{g}_2
    $$
\end{lemme}

\begin{proposition}{}{}
    Soit $f \in Cpm([0, 2\pi[, \mathbb{C})$. Alors
    $$
    |\hat{f_k}| \leq \norm{f}_2
    $$
    pour tout $k \in \mathbb{Z}$
\end{proposition}

\begin{demonstration}
    En a bien
    $$
    \hat{f_k} = \langle f, e^{ikt} \rangle \leq \norm{f}_2 \norm{e^{ikt}}_2 = \norm{f}_2
    $$
\end{demonstration}

\begin{definition}
    Soient $f, g \in Cpm([0, 2\pi[, \mathbb{C})$.
    
    $f$ et $g$ sont dits orthogonaux si $\langle f, g \rangle = 0$.
\end{definition}

\begin{theoreme}{}{}
    Pour tout $k \in \Z$ si on note $e_k = e^{-ikx}$, alors la famille $(e_k)_{k \in \Z}$ est orthogonale dans l'espace des fonctions intégrables sur $[0, 2\pi[$.
\end{theoreme}

\begin{demonstration}
    Immédiatement, on a
    $$
    \langle e_k, e_j \rangle = \frac{1}{2\pi} \int_0^{2\pi} e_k(t) \overline{e_j(t)} \, dt = \frac{1}{2\pi} \int_0^{2\pi} e^{i(j-k)t} \, dt = \begin{cases}
        1 & \text{si } j = k \\
        0 & \text{sinon}
    \end{cases} 
    $$    
\end{demonstration}

\begin{corollaire}{}{}
    La famille $\{ \cos(kx), \sin(kx) \}_{k \in \N}$ est orthogonale dans l'espace des fonctions intégrables sur $[0, 2\pi[$.
\end{corollaire}

\begin{corollaire}{}{}
    Si $\tau$ est un polynôme trigonométrique d'ordre $n$ de la forme
    $$
    \tau(x) = \sum_{j=-n}^n c_j (e^{ix})^j
    $$
    alors $c_j = \hat{\tau_j}$ pour tout $j \in \{-n, \dots, n\}$.
\end{corollaire}

\begin{demonstration}
    On a
    $$
    \hat{\tau_k} = \frac{1}{2\pi} \int_0^{2\pi} \tau(t) e^{-ikt} \, dt = \frac{1}{2\pi} \int_0^{2\pi} \sum_{j=-n}^n c_j e^{ijt} e^{-ikt} \, dt = \sum_{j=-n}^n c_j \frac{1}{2\pi} \int_0^{2\pi} e^{i(j-k)t} \, dt = c_k
    $$
\end{demonstration}

\begin{corollaire}{}{}
    Si $\tau$ est un polynôme trigonométrique d'ordre $n$, alors
    $$
    \norm{\tau}_2^2 = \sum_{j=-n}^n |\hat{\tau_j}|^2
    $$
\end{corollaire}

\begin{demonstration}
    On a
    \begin{align*}
        \norm{\tau}_2^2 = \langle \tau, \tau \rangle &= \left\langle \sum_{j=-n}^n \hat{\tau_j} e^{ijt}, \sum_{k=-n}^n \hat{\tau_k} e^{ikt} \right\rangle \\
        &= \frac{1}{2\pi} \int_0^{2\pi} \left( \sum_{j=-n}^n \hat{\tau_j} e^{ijt} \right) \left( \sum_{k=-n}^n \overline{\hat{\tau_k}} e^{-ikt} \right) \, dt \\
        &= \sum_{j=-n}^n \sum_{k=-n}^n \hat{\tau_j} \overline{\hat{\tau_k}} \frac{1}{2\pi} \int_0^{2\pi} e^{i(j-k)t} \, dt \\
        &= \sum_{j=-n}^n \hat{\tau_j} \overline{\hat{\tau_j}} \\
        &= \sum_{j=-n}^n |\hat{\tau_j}|^2
    \end{align*}
\end{demonstration}

\begin{theoreme}{}{}
    Soit $f \in Cpm([0, 2\pi[, \mathbb{C})$ bornée.

    $S_{f_n}$ la somme partielle d'ordre $n$ de la série de Fourier de $f$.
    $$
    S_{f_n}(t) = \sum_{j=-n}^n \hat{f_j} e^{ijt}
    $$
    et soit $\tau$ un polynôme trigonométrique d'ordre $n$. Alors
    $$
    \norm{f - S_{f_n}}_2 \leq \norm{f - \tau}_2
    $$
\end{theoreme}

\begin{demonstration}
    On a 
    $$
    \norm{f - \tau}_2^2 = \langle f - \tau, f - \tau \rangle = \norm{f}_2^2 - \langle f, \tau \rangle - \overline{\langle f, \tau \rangle} + \norm{\tau}_2^2 \\
    $$
    on injecte
    \begin{align*}
        \langle f, \tau \rangle &= \frac{1}{2\pi} \int_0^{2\pi} f(t) \overline{\tau(t)} \, dt \\
        &= \frac{1}{2\pi} \int_0^{2\pi} f(t) \sum_{j=-n}^n \overline{\hat{\tau_j}} e^{-ijt} \, dt \\
        &= \sum_{j=-n}^n \overline{\hat{\tau_j}} \frac{1}{2\pi} \int_0^{2\pi} f(t) e^{-ijt} \, dt \\
        &= \langle f, S_{f_n} \rangle
    \end{align*}
    et donc
    \begin{align*}  
        \norm{f - \tau}_2^2 &= \norm{f}_2^2 - \norm{S_n}_2^2 + \norm{S_n}_2^2 - \langle f, S_{f_n} \rangle - \overline{\langle f, S_{f_n} \rangle} + \norm{S_{f_n}}_2^2 \\
        &= \norm{f}_2^2 - \norm{S_n}_2^2 + \norm{S_n - \tau}_2^2 \\
        &\geq \norm{f}_2^2 - \norm{S_n}_2^2
    \end{align*}

    Si on considère $\tau = S_{f_n}$, on a
    $$
    \norm{f - S_{f_n}}_2^2 = \norm{f}_2^2 - \langle f, S_{f_n} \rangle - \overline{\langle f, S_{f_n} \rangle} + \norm{S_{f_n}}_2^2
    $$
    et
    $$
    \langle f, S_{f_n} \rangle = \langle S_{f_n}, S_{f_n} \rangle = \norm{S_{f_n}}_2^2
    $$
    donc
    $$
    \norm{f - S_{f_n}}_2^2 = \norm{f}_2^2 - \norm{S_{f_n}}_2^2 - \norm{S_{f_n}}_2^2 + \norm{S_{f_n}}_2^2 = \norm{f}_2^2 - \norm{S_{f_n}}_2^2
    $$
\end{demonstration}

\begin{corollaire}{Inégalité de Bessel}{}
    Soit $f: [0, 2\pi[ \to \mathbb{C}$ bornée et intégrable sur $[0, 2\pi[$. Alors
    $$
    \sum_{j=-\infty}^{+\infty} |\hat{f_j}|^2 \leq \norm{f}_2^2
    $$
    et en particulier
    $$
    \lim_{k \to +\infty} |\hat{f_k}| = 0
    $$
\end{corollaire}

\begin{demonstration}
    On sait déjà que $\norm{f - S_{f,n}}_2^2 = \norm{f}_2^2 - \norm{S_{f,n}}_2^2 \geq 0$
    donc $\norm{S_{f,n}}_2^2 \leq \norm{f}_2^2$ et de plus $\norm{S_{f,n}}_2^2 = \sum_{j=-n}^n |\hat{f_j}|^2$.

    Dans toutes les sommes partielles $S_{f,n}$ sont bornée et $S_{f,n}$ est une série à terme positifs,
    donc la série est convergente et avec $\sum\limits_{j=-\infty}^{+\infty} |\hat{f_j}|^2 \leq \norm{f}_2^2$, on a $\lim\limits_{k \to +\infty} |\hat{f_k}|^2 = 0$ et donc $\lim\limits_{k \to +\infty} |\hat{f_k}| = 0$.
\end{demonstration}

\begin{corollaire}{Identité de Parseval}{}
    Soit $f: [0, 2\pi[ \to \mathbb{C}$ bornée et intégrable sur $[0, 2\pi[$.

    Si pour tout $\varepsilon > 0$, il existe un $N_\varepsilon \geq 0$ et un polynôme trigonométrique d'ordre $N_\varepsilon$ pour lequel $\max\limits_{x \in [0, 2\pi[} |f(x) - \tau(x)| \leq \varepsilon$.

    Alors on a
    $$
    \sum_{j=-\infty}^{+\infty} |\hat{f_j}|^2 = \norm{f}_2^2
    $$
\end{corollaire}

\begin{demonstration}
    Soit $\varepsilon > 0$ donné, $N_\varepsilon$ et $\tau$ existent
    \begin{align*}
        0 \leq \norm{f - S_{f,N_\varepsilon}}_2^2 &= \norm{f}_2^2 - \norm{S_{f,N_\varepsilon}}_2^2 \\
        &\leq \norm{f - \tau}_2^2 \\
        & \leq \frac{1}{2\pi} \int_0^{2\pi} |f(t) - \tau(t)|^2 \, dt \\
        &\leq \varepsilon^2
    \end{align*}

    or $\norm{S_{f,N_\varepsilon}}_2^2 = \sum_{j=-N_\varepsilon}^{N_\varepsilon} |\hat{f_j}|^2$, cette série est àtp, croissante et convergente.

    ainsi pour tout $n \geq N_\varepsilon$, on a
    $$
    0 \leq \norm{f}_2^2 - \sum_{j=-n}^n |\hat{f_j}|^2 \leq \norm{f}_2^2 - \sum_{j=-N_\varepsilon}^{N_\varepsilon} |\hat{f_j}|^2 \leq \varepsilon^2
    $$

    On peut passer à la limite quand $n \to +\infty$ et on a le résultat.
\end{demonstration}

\subsection{Approximations ponctuelles et uniformes}

\begin{definition}[(Le noyau de Dirichlet)]
    On appelle \textbf{noyau de Dirichlet} d'ordre $k$, noté $D_k$, le polynôme trigonométrique d'ordre $k$ dont tous les coefficients sont égaux à $\frac{1}{2\pi}$, c'est à dire
    $$
    D_k(t) = \frac{1}{2\pi} \sum_{j=-k}^k e^{-ijt}
    $$
\end{definition}

\begin{proposition}{}{}
    $$
    D_k(t) = \frac{1}{2\pi} \frac{\sin((k + \frac{1}{2})t)}{\sin(\frac{t}{2})} \quad \text{et} \quad \int_{-\pi}^{\pi} D_k(t) \, dt = 1
    $$
    et $D_k$ est une fonction paire.
    \begin{hotwarn}
        exercice maison
    \end{hotwarn}
\end{proposition}

\begin{definition}[(Le noyau de Fejér)]
    On appelle \textbf{noyau de Fejér} d'ordre $k$, noté $K_n$, la moyenne arithmétique des noyaux de Dirichlet d'ordre inférieur ou égal à $n$, c'est à dire
    $$
    K_n(t) = \frac{1}{n + 1} \sum_{k=0}^n D_k(t)
    $$
\end{definition}

\begin{proposition}{}{}
    $$
    K_n(t) = \frac{\sin^2(\frac{(n + 1)t}{2})}{2\pi (n + 1) \sin^2(\frac{t}{2})} \quad \text{et} \quad \int_{-\pi}^{\pi} K_n(t) \, dt = 1
    $$
    \begin{hotwarn}
        exercice maison
    \end{hotwarn}
\end{proposition}

\end{document}